\section{Kajian Pustaka}

\subsection{Peto's Paradox}

Peto's Paradox berangkat dari ketidaksesuaian antara prediksi sederhana dan
data lintas spesies. Jika setiap sel memiliki peluang yang sama untuk
mengalami transformasi kanker, maka organisme besar dan berumur panjang
diperkirakan memiliki risiko kanker lebih tinggi daripada organisme
kecil. Pada kenyataannya, korelasi tersebut lemah atau tidak proporsional
antarspesies \cite{caulin2011}. Hal ini berbeda dengan pola dalam satu
spesies, di mana ukuran tubuh atau umur sering berkaitan dengan risiko
kanker yang lebih tinggi.

Penjelasan biologis yang umum diajukan adalah evolusi mekanisme supresi
kanker. Mekanisme tersebut dapat berupa DNA repair yang lebih efisien,
checkpoint yang lebih ketat, apoptosis yang lebih sensitif, pembatasan
pembelahan stem cell, arsitektur jaringan yang membatasi ekspansi klon,
atau sistem imun yang lebih efektif. Contoh yang banyak dibahas adalah
gajah. Sulak dkk. melaporkan ekspansi salinan TP53 pada gajah dan respons
kerusakan DNA yang lebih aktif \cite{sulak2016}. Tollis dkk. juga melaporkan
percepatan evolusi pada mekanisme pertahanan penyakit gajah
\cite{tollis2021}.

\subsection{Dataset Kanker Lintas Mamalia}

Vincze dkk. menyusun basis data cancer-related mortality pada mamalia dewasa
di bawah perawatan manusia \cite{vincze2022}. Data berasal dari Species360
dan Zoological Information Management System. Analisis utama mereka
melibatkan 191 spesies mamalia dan menunjukkan bahwa cancer mortality risk
sebagian besar independen dari massa tubuh dan adult life expectancy.
Mereka juga menekankan perbedaan antarordo dan kaitan diet dengan risiko
kanker.

Dalam paper ini, data yang tersedia publik dari GitHub penulis digunakan
sebagai input analisis. Kolom utama yang dipakai adalah Species, Order,
Species body mass, Adult life expectancy, CMR, ICM, jumlah individu, jumlah
catatan postmortem, dan jumlah kasus neoplasia. Setelah baris dengan nilai
inti kosong dikeluarkan, analisis menghasilkan 172 spesies. Dari jumlah
tersebut, 131 spesies memiliki CMR lebih dari nol.

\subsection{Siklus Sel sebagai Interpretasi Mekanistik}

Data komparatif tidak mengukur DNA repair, checkpoint, atau apoptosis secara
langsung pada setiap spesies. Karena itu, paper ini tidak mengklaim bahwa
residual statistik adalah indikator langsung satu mekanisme molekuler tertentu.
Sebaliknya, residual dipakai sebagai sinyal komputasional: jika suatu
spesies memiliki exposure proxy tinggi tetapi CMR rendah, maka spesies
tersebut menjadi kandidat untuk diteliti mekanisme supresinya.

Interpretasi siklus sel tetap relevan karena sebagian besar mekanisme
supresi kanker bekerja dengan mengurangi peluang mutasi bertahan dan
berekspansi. DNA repair menurunkan mutasi efektif. Checkpoint menghambat
kelanjutan siklus sel saat kerusakan terdeteksi. Apoptosis menghilangkan
sel berisiko. Dengan demikian, analisis komparatif dapat digunakan untuk
menentukan spesies atau ordo yang layak diprioritaskan dalam studi mekanisme
kontrol siklus sel.
